% Part: first-order-logic
% Chapter: syntax-and-semantics
% Section: introduction

\documentclass[../../../include/open-logic-section]{subfiles}

\begin{document}

\olfileid{fol}{syn}{int}

\olsection{పరిచయం}

మొదటిస్థాయి తర్కశాస్త్రపు సిద్ధాంతం, అధిసిద్ధాంతం అభివృద్ధి చేయాలంటే ముందుగా దాని వ్యక్తీకరణల వాక్యనిర్మాణం, అర్థవిచారాన్ని నిర్వచించాలి. మొదటిస్థాయి తర్కశాస్త్రపు వ్యక్తీకరణలు పదాలు, \tetoken{సూత్రాలు}{!!{formula}s}. పదాలు \tetoken{చరరాశులు}{!!{variable}s}, \tetoken{స్థిరాంకాలు}{!!{constant}s}, \tetoken{ప్రమేయ సంకేతాలతో}{!!{function}s} ఏర్పడతాయి. \tetoken{సూత్రాలు}{!!^{formula}s} అయితే \tetoken{విధేయ సంకేతాలు}{!!{predicate}s}, పదాలతో ఏర్పడతాయి (ఇవే అతి చిన్న, ``అణు'' \tetoken{సూత్రాలు}{!!{formula}s}); తరువాత తార్కిక సంయోజకాలు, పరిమాణీకరణికాలను ఉపయోగించి అణు \tetoken{సూత్రాల}{!!{formula}s} నుంచి మరింత సంక్లిష్ట సూత్రాలను నిర్మించవచ్చు. నిర్మాణ నియమాలను పేర్కొనడానికి ఎన్నో మార్గాలున్నాయి; వాటిలో సాధ్యమైన ఒక విధానాన్నే ఇక్కడ ఇస్తాం. ఇతర వ్యవస్థలు వేరే సంకేతాలను ఎంచుకోవచ్చు, వేరే సంయోజకాల సమూహాలను మూలభూతాలుగా తీసుకోవచ్చు, కుండలీకరణాలను వేరుగా వాడవచ్చు (లేదా పోలిష్ సంకేతనం మాదిరిగా అసలు వాడకపోవచ్చు). అన్ని విధానాలకూ ఉమ్మడిగా ఉండేదేమిటంటే, నిర్మాణ నియమాలు పదాలూ \tetoken{సూత్రాలూ}{!!{formula}s} ఉన్న సమితిని \emph{ఆగమనాత్మకంగా} నిర్వచిస్తాయి. సరిగ్గా చేస్తే, ప్రతి వ్యక్తీకరణ నిర్మాణ నియమాల ప్రకారం ప్రధానంగా ఒకే విధంగా ఏర్పడుతుంది. \emph{ఏకైక పఠనయోగ్యత} గల వ్యక్తీకరణలను ఇచ్చే ఈ ఆగమనాత్మక నిర్వచనం వల్ల, అదే పద్ధతిలో---ఆగమనాత్మక నిర్వచనంతోనే---వాటికి అర్థాలను ఇవ్వగలం.

వ్యక్తీకరణలకు అర్థం ఇవ్వడం అర్థవిచార పరిధిలోని పని. అర్థవిచారంలో కేంద్ర భావన \tetoken{ఒక నిర్మాణంలో}{!!a{structure}} సంతృప్తి చెందడం. \tetoken{ఒక నిర్మాణం}{!!^a{structure}} భాషలోని మూల భాగాలకు అర్థాన్ని ఇస్తుంది: \tetoken{ఒక ప్రదేశం}{!!a{domain}} అనేది వస్తువుల ఖాళీ కాని సమితి. పరిమాణీకరణికాలు ఈ ప్రదేశంపై విస్తరించేవిగా అన్వయించబడతాయి; \tetoken{స్థిరాంకాలకు}{!!{constant}s} ప్రదేశంలోని మూలకాలు కేటాయించబడతాయి; \tetoken{ప్రమేయ సంకేతాలకు}{!!{function}s} \tetoken{ప్రదేశం}{!!{domain}} నుంచి దానికే ప్రమేయాలు కేటాయించబడతాయి; \tetoken{విధేయ సంకేతాలకు}{!!{predicate}s} \tetoken{ప్రదేశం}{!!{domain}}పై సంబంధాలు కేటాయించబడతాయి. \tetoken{ప్రదేశం}{!!{domain}}తో పాటు మూల పదజాలానికి ఇచ్చిన కేటాయింపులు కలిపి \tetoken{ఒక నిర్మాణం}{!!a{structure}} అవుతుంది. \tetoken{చరరాశులు}{!!^{variable}s}, \tetoken{సూత్రాలలో}{!!{formula}s} ఉండవచ్చు; వాటికి అర్థవిచారం ఇవ్వాలంటే \tetoken{ప్రదేశంలోని}{!!{domain}} \tetoken{మూలకాలను}{!!{element}s} కూడా వాటికి కేటాయించాలి---దానినే చరరాశి కేటాయింపు అంటాం. చివరగా సంతృప్తి సంబంధం వీటన్నింటినీ కలుపుతుంది. ఒక \tetoken{సూత్రం}{!!^a{formula}}, చరరాశి కేటాయింపు~$s$కు సాపేక్షంగా \tetoken{ఒక నిర్మాణం}{!!a{structure}}~$\Struct{M}$లో సంతృప్తి చెందవచ్చు; దాన్ని $\Sat{M}{!A}[s]$ అని రాస్తాం. ఈ సంబంధాన్ని కూడా $!A$ నిర్మాణంపై ఆగమన పద్ధతిలో నిర్వచిస్తాం; ఉదాహరణకు $!A \land !B$ సంతృప్తిని $!A$, ~$!B$ సంతృప్తి (లేదా అసంతృప్తి) పరంగా చెప్పడానికి తార్కిక సంయోజకాల సత్య పట్టికలను వాడతాం. అప్పుడు \tetoken{సూత్రం}{!!{formula}}~$!A$ ఒక \tetoken{వాక్యం}{!!a{sentence}} అయితే, అంటే స్వేచ్ఛా చరరాశులు లేకపోతే, చరరాశి కేటాయింపు సంబంధంలేనిదని తెలుస్తుంది. కాబట్టి \tetoken{వాక్యాలు}{!!{sentence}s} \tetoken{నిర్మాణాలలో}{!!{structure}s} సంతృప్తి చెందుతాయి (లేదా చెందవు) అని నేరుగా చెప్పవచ్చు.

వాక్యాల కోసం సంతృప్తి సంబంధం $\Sat{M}{!A}$ ఆధారంగా, చెల్లుబాటు, అనుగమనం, సంతృప్త్యర్హత అనే ప్రాథమిక అర్థవిచార భావనలను నిర్వచించవచ్చు. ప్రతి నిర్మాణం సంతృప్తి పరచే వాక్యం చెల్లుబాటైనది, $\Entails !A$. ఒక \tetoken{వాక్యాల}{!!{sentence}s} సమితి నుంచి వాక్యం అనుగమిస్తుంది, $\Gamma \Entails !A$, అంటే~$\Gamma$లోని అన్ని \tetoken{వాక్యాలను}{!!{sentence}s} సంతృప్తి పరచే ప్రతి \tetoken{నిర్మాణం}{!!{structure}}~$!A$ను కూడా సంతృప్తి పరుస్తుంది. అలాగే ఏదో ఒక \tetoken{నిర్మాణం}{!!{structure}} ఒక సమితిలోని అన్ని \tetoken{వాక్యాలను}{!!{sentence}s} ఒకేసారి సంతృప్తి పరిస్తే ఆ వాక్యాల సమితి సంతృప్త్యర్హమైనది. \tetoken{సూత్రాలు}{!!{formula}s} ఆగమనాత్మకంగా నిర్వచించబడటం, సంతృప్తిని మళ్లీ \tetoken{సూత్రాల}{!!{formula}s} నిర్మాణంపై ఆగమన పద్ధతిలో నిర్వచించడం వల్ల, అర్థవిచారపు ధర్మాలను నిరూపించడానికి, నిర్వచించిన అర్థవిచార భావనలను పరస్పరం అనుసంధానించడానికి ఆగమన పద్ధతిని ఉపయోగించవచ్చు.




\end{document}
